\documentclass[11pt]{article}
\usepackage[margin=1in]{geometry}
\usepackage[T1]{fontenc}
\usepackage{lmodern}
\usepackage{amsmath,amssymb,mathtools,bm}
\usepackage{booktabs,tabularx,array,longtable}
\usepackage{graphicx}
\usepackage{xcolor}
\usepackage{float}
\usepackage{placeins}
\usepackage{microtype}
\usepackage{fancyhdr}
\usepackage{hyperref}
\usepackage{xurl}
\hypersetup{hidelinks,pdfdisplaydoctitle=true,
 pdftitle={HPS Gaussian-Process Resonance Search: Signal-Tail Study v5.9},
 pdfauthor={Emrys Peets},
 pdfsubject={Observed signal-tail scenarios for full 2015, full 2016 and 2021 10 percent}}
\urlstyle{same}
\pagestyle{fancy}
\fancyhf{}
\setlength{\headheight}{14pt}
\fancyhead[L]{\small HPS Gaussian-Process Resonance Search}
\fancyhead[R]{\small v5.9 / Signal-tail study}
\fancyfoot[C]{\thepage}
\fancypagestyle{plain}{\fancyhf{}
 \fancyhead[L]{\small HPS Gaussian-Process Resonance Search}
 \fancyhead[R]{\small v5.9 / Signal-tail study}
 \fancyfoot[C]{\thepage}}
\setlength{\tabcolsep}{5pt}
\setlength{\emergencystretch}{2em}
\renewcommand{\arraystretch}{1.12}
\newcolumntype{Y}{>{\raggedright\arraybackslash}X}
\newcommand{\eps}{\varepsilon}
\newcommand{\CLs}{CL$_{\mathrm{s}}$}
\newcommand{\sig}{\sigma_m}
\newcommand{\studfig}[4]{\begin{figure}[H]\centering
 \includegraphics[width=#1\linewidth,height=#2\textheight,keepaspectratio]{#3}
 \caption{#4}\end{figure}}
\newcommand{\studtable}[2]{\begin{table}[H]\centering\small
 \input{#1}\caption{#2}\end{table}}
\title{HPS Gaussian-Process Resonance Search\\[0.5em]
 \large Signal-tail study v5.9\\
 Gaussian cores with broader tails}
\author{Emrys Peets\\Stanford University and SLAC National Accelerator Laboratory}
\date{21 September 2026\\\normalsize Standalone study following Analysis Note v5.0.5}

\begin{document}
\maketitle
\begin{abstract}
This study tests symmetric signal shapes with a Gaussian core inside the
primary $\pm2.25\sig$ blind window and broader tails outside it. The tail
scale is increased by 10, 20 and 30\% in two specified shape families.
Observed 90\% \CLs{} limits and fixed-mass, one-sided asymptotic
background-only probabilities are evaluated for full 2015 over 19--100~MeV,
full 2016 over 39--180~MeV, and the released 2021 10\% sample over
50--250~MeV. Each family is compared with a Gaussian under the same
background and window prescription. The primary calculation retains the
v5.0.5 window; a separate $\pm4.5\sig$ calculation tests the information
carried by the outer structure. Detector resolution, prompt-density
normalization and archived kernel coordinates are held fixed. The results
retain primary-window maxima at 51, 90 and 78~MeV, respectively. The
smooth-dilation alternatives increase the median primary limits by
approximately 0.168, 0.345 and 0.529\%, with nearly unchanged local
probabilities. These are declared signal-shape scenarios, conditional on the stated choices.
They do not provide coverage calibration or a new global discovery probability.
\end{abstract}

\section{Observed comparisons and scope}
The requested tail changes are implemented outside a fixed Gaussian core.
They are changes in a tail-scale parameter, rather than a 10--30\% increase
in the detector resolution or in the total number of signal events. Each
template is normalized to unit yield over the inherited support, and the
fitted amplitude denotes that full-template yield.

There are seven shapes in each window: one Gaussian, three smooth tail
dilations and three tail-curvature variations. The three datasets are
analyzed separately. No dataset combination or choice among scenarios is
used to strengthen an observed limit or significance. The complete
mass-by-mass results accompany the note; the peak tables identify the
smallest local probability within each stated scan, without assigning a
probability to the selection of that minimum.

\clearpage
\subsection{The primary-window results}
Table~\ref{tab:primary-peaks} records the strongest local excess for every
primary-window scenario. The Gaussian row is freshly computed with the
same portable fitting machinery used for all alternatives. For each
alternative, the coupling and yield limits refer to its
own full-template normalization, while the local probability tests whether
a positive signal component improves the fit at that mass.

The primary-window local maxima remain at 51~MeV (2015), 90~MeV (2016), and 78~MeV (2021 10\%) for all six altered shapes. For the primary dilation model, the median observed limit increases by approximately 0.168\%, 0.345\%, and 0.529\% for 10\%, 20\%, and 30\% broader tails. These small shifts follow the change in core probability; they are not 10--30\% changes in the full-template resolution.

In the guard window, the 30\% dilation changes observed limits by $-0.048$ to $+2.068$\% in 2015, $+0.245$ to $+2.056$\% in 2016, and $+0.279$ to $+1.543$\% in 2021. The largest absolute changes in the signed local root are 0.0377, 0.0280, and 0.0177, respectively. The guard-window Gaussian maximum in 2015 is already at 93~MeV, rather than 51~MeV. That movement is a consequence of the wider background interpolation and fitted region; it cannot be attributed to the tail variation. The primary search is complete over all declared mass intervals. The guard search omits 2016 masses 176--180~MeV because fewer than three upper exterior bins remain.


\begin{table}[H]\centering\small
\begin{tabular}{llrrrr}
\toprule
Dataset & Shape & $m_0$ (MeV) & $Z_0$ & Local $p_0$ & $\varepsilon^2_{90}$ \\
\midrule
2015 & Gaussian & 51 & 3.13924 & $8.469\times10^{-4}$ & $2.640\times10^{-5}$ \\
2015 & Dilation +10\% & 51 & 3.13925 & $8.469\times10^{-4}$ & $2.644\times10^{-5}$ \\
2015 & Dilation +20\% & 51 & 3.13926 & $8.469\times10^{-4}$ & $2.649\times10^{-5}$ \\
2015 & Dilation +30\% & 51 & 3.13927 & $8.469\times10^{-4}$ & $2.654\times10^{-5}$ \\
2015 & Curvature +10\% & 51 & 3.13924 & $8.469\times10^{-4}$ & $2.641\times10^{-5}$ \\
2015 & Curvature +20\% & 51 & 3.13924 & $8.469\times10^{-4}$ & $2.642\times10^{-5}$ \\
2015 & Curvature +30\% & 51 & 3.13924 & $8.469\times10^{-4}$ & $2.643\times10^{-5}$ \\
2016 & Gaussian & 90 & 3.42475 & $3.077\times10^{-4}$ & $1.313\times10^{-5}$ \\
2016 & Dilation +10\% & 90 & 3.42475 & $3.077\times10^{-4}$ & $1.316\times10^{-5}$ \\
2016 & Dilation +20\% & 90 & 3.42475 & $3.077\times10^{-4}$ & $1.318\times10^{-5}$ \\
2016 & Dilation +30\% & 90 & 3.42475 & $3.077\times10^{-4}$ & $1.320\times10^{-5}$ \\
2016 & Curvature +10\% & 90 & 3.42475 & $3.077\times10^{-4}$ & $1.314\times10^{-5}$ \\
2016 & Curvature +20\% & 90 & 3.42475 & $3.077\times10^{-4}$ & $1.315\times10^{-5}$ \\
2016 & Curvature +30\% & 90 & 3.42475 & $3.077\times10^{-4}$ & $1.315\times10^{-5}$ \\
2021 & Gaussian & 78 & 2.80864 & $2.488\times10^{-3}$ & $6.330\times10^{-6}$ \\
2021 & Dilation +10\% & 78 & 2.80864 & $2.488\times10^{-3}$ & $6.341\times10^{-6}$ \\
2021 & Dilation +20\% & 78 & 2.80864 & $2.488\times10^{-3}$ & $6.352\times10^{-6}$ \\
2021 & Dilation +30\% & 78 & 2.80864 & $2.488\times10^{-3}$ & $6.363\times10^{-6}$ \\
2021 & Curvature +10\% & 78 & 2.80864 & $2.488\times10^{-3}$ & $6.333\times10^{-6}$ \\
2021 & Curvature +20\% & 78 & 2.80864 & $2.488\times10^{-3}$ & $6.336\times10^{-6}$ \\
2021 & Curvature +30\% & 78 & 2.80864 & $2.488\times10^{-3}$ & $6.338\times10^{-6}$ \\
\bottomrule
\end{tabular}

\caption{Primary-window scan minima and corresponding observed limits.
Every probability is local, with the mass and shape fixed. A scan minimum
does not include a mass or scenario trials correction.}
\label{tab:primary-peaks}
\end{table}

\clearpage
\section{Signal shapes and normalization}
Let $z=|m-m_0|/\sig$, $a=2.25$ and $t=z-a$. The nominal resolution
$\sig=\sig(m_0)$ is inherited from v5.0.5. All shapes have the
unnormalized core
\begin{equation}
 f(z)=\exp(-z^2/2),\qquad 0\leq z\leq a.
\end{equation}
For $z>a$, the primary family stretches the tail coordinate smoothly:
\begin{align}
 u_k(z)&=a+\frac{t}{1+(k-1)[1-\exp(-t/h)]},\qquad h=0.5,\\
 f_{\mathrm{dil},k}(z)&=\exp[-u_k(z)^2/2],
 \qquad k\in\{1.1,1.2,1.3\}.
 \label{eq:dilation}
\end{align}
The coordinate and first derivative match the Gaussian at $z=a$;
far into the tail, the distance from the boundary is stretched by $k$.
The secondary family instead changes the tail curvature:
\begin{equation}
 f_{\mathrm{curv},k}(z)=
 \exp\!\left[-\frac{a^2}{2}-at-\frac{t^2}{2k^2}\right],
 \qquad z>a.
 \label{eq:curvature}
\end{equation}
It also matches the core value and slope at the boundary. The two
families test different continuations of the same core; their common $k$
does not imply equal tail fractions or equal root-mean-square widths.
Both reduce exactly to the Gaussian at $k=1$.

\studfig{0.98}{0.34}{../figures/shapes_note.pdf}{Smooth-dilation signal
shapes and the 30\% curvature control in units of the nominal mass
resolution. The left panel shows full-line-normalized densities and the
right shows their ratios to the Gaussian. Changes to the unnormalized
shape begin at the primary blind-window boundary.}

Core-only bins differ solely by normalization, leaving their unrestricted
likelihood ratio and $p_0$ invariant. The actual center-selected bins can
straddle $\pm2.25\sig$; their integrated tails break this proportionality
and can produce small primary-probability changes. The guard calculation
explicitly tests outer bins.

\clearpage
\subsection{Bin integration and the full-template yield}
For support $\mathcal D_d$ and analysis bin $B_i$ in dataset $d$, the
template probability is
\begin{equation}
 T_{di}^{(s)}(m_0)=
 \frac{\displaystyle\int_{B_i}f_s(|m-m_0|/\sig)\,dm}
 {\displaystyle\int_{\mathcal D_d}f_s(|m-m_0|/\sig)\,dm},
 \qquad \sum_{i\in\mathcal D_d}T_{di}^{(s)}=1.
 \label{eq:bin-template}
\end{equation}
The integrated bin probabilities are retained when a fit window is
selected. They are never renormalized inside that window. If
$F_{d,s,W}=\sum_{i\in W}T_{di}^{(s)}$, the fitted-window yield at the
upper limit is $A_{90}F_{d,s,W}$, while $A_{90}$ itself denotes the
full-template yield. The normalization in Eq.~\eqref{eq:bin-template}
also accounts for the finite support near its endpoints.

\studtable{../derived/shape_table.tex}{Numerical shape properties for
the declared tail prescriptions, evaluated on the full line. The core
upper-limit ratio is the ideal proportional-core normalization ratio.
Actual fitted templates use the finite support in
Eq.~\eqref{eq:bin-template}.}

The prompt-density normalization uses the same nominal resolution for
every shape:
\begin{equation}
 \rho_d(m_0)=\frac{1}{3.28\sig}
 \int_{m_0-1.64\sig}^{m_0+1.64\sig}\frac{dN_d}{dm}\,dm,
 \qquad
 K_d(m_0)=\frac{3\pi m_0 f_{\mathrm{rad},d}^{\mathrm{eff}}}
 {2\alpha}\,\rho_d(m_0).
 \label{eq:conversion}
\end{equation}
The integral uses fractional overlaps with the native histogram bins.
The electron-channel coupling coordinate satisfies $A=K_d\eps^2$.
Holding $K_d$ fixed isolates the specified shape change; this study does
not infer a new efficiency or acceptance from the broader tails.

Above the dimuon threshold, the inherited minimal-visible coupling
correction is applied once:
\begin{equation}
 (\eps^2_{90})_{\mathrm{minimal}}=\frac{A_{90}}{K_d}\,c_\mu(m_0),
 \quad c_\mu(m_0)=
 \begin{cases}
 1,&m_0\leq2m_\mu,\\
 1+\sqrt{1-4r_\mu}(1+2r_\mu),&m_0>2m_\mu,
 \end{cases}
 \quad r_\mu=(m_\mu/m_0)^2.
\end{equation}
This conversion changes neither the visible signal-yield limit nor $p_0$.
In the numerical ledger, \texttt{A90} is the solver coordinate
$\eps^2/(10^{-8})$, whereas \texttt{full\_yield\_90} is the event yield
denoted $A_{90}$ in this note; \texttt{derived/COLUMN\_GUIDE.md} defines
the columns and conversions.

\clearpage
\section{Primary-window observed scans}
\subsection{Full 2015: 19--100 MeV}
The 2015 scan retains the 14--135~MeV support and 0.25~MeV analysis bins.
The 91--100~MeV extension follows the v5.0.5 endpoint prescription:
kernel coordinates are fixed at their archived 90~MeV values, while the
resolution, mask, template and density are evaluated at the actual mass.
The extrapolated interval keeps its conditional status.

\studfig{0.98}{0.69}{../figures/primary_2015.pdf}{Observed 2015
primary-window Gaussian and smooth-dilation scenarios. The coupling
conversion refers to the complete normalized template; the likelihood uses the selected
$\pm2.25\sig$ bins. Limits are pointwise asymptotic 90\% \CLs{} limits
and probabilities are fixed-mass local references.}

\clearpage
\subsection{Full 2016: 39--180 MeV}
The 2016 scan uses the frozen 30--210~MeV support and 0.25~MeV analysis
bins. Every shape at a mass uses the same archived kernel coordinates
and the same background constraint for this window. Thus the
within-window comparison changes the signal template without selecting
a new background fit from the resulting limit or probability.

The v5.0.5 note documents differences in independent 2016
hyperparameter fits. Reusing its archived states makes this calculation
reproducible under that prescription; it does not settle the separate
fit-provenance or background-calibration questions.

\studfig{0.98}{0.65}{../figures/primary_2016.pdf}{Observed 2016
primary-window Gaussian and smooth-dilation results. The same frozen-state qualification applies to
the Gaussian and every tail alternative. Probability differences should
be interpreted together with the small boundary-bin contributions and
the explicitly broader calculation below.}

\clearpage
\subsection{Released 2021 10\%: 50--250 MeV}
The 2021 scan uses the native 10\% sample with the inherited
36--300~MeV support prescription. Rebinning leaves the final complete
analysis-bin edge at 299.75~MeV; the analysis-bin width is 0.625~MeV.
The larger bin width makes the distinction between a continuous mask
boundary and a full selected bin particularly relevant to the core-only
comparison.

\studfig{0.98}{0.69}{../figures/primary_2021.pdf}{Observed 2021 10\%
primary-window Gaussian and smooth-dilation results. Above $2m_\mu$, the plotted coupling includes the
inherited branching correction once. Neither the visible yield nor the
local probability receives this correction.}

\clearpage
\section{Outer structure in a wider fit region}
The guard calculation expands both the fitted region and the training
exclusion to $\pm4.5\sig$. Its GP is conditioned only on bins outside
that wider mask, using the same archived kernel coordinates. The
prediction and covariance are recomputed for the widened region. Every
observed bin in that region enters the likelihood once.

The primary scans contain all 425 requested dataset--mass coordinates.
The guard scan contains 420: the five 2016 hypotheses at 176--180~MeV
do not retain the required external sidebands under the wider mask and
are excluded by the geometric rule. Its 2016 domain is consequently
39--175~MeV. The 2015 and 2021 guard domains retain their full primary
ranges. Across seven shapes and both windows, 5,915 observed fits are
reported; omitted guard endpoints are not interpolated or filled from the
primary calculation.

This is a separate background-interpolation and extraction prescription.
The meaningful shape comparison is therefore the tail scenario divided
by, or compared with, the Gaussian \emph{in the same guard window}.
Comparing a guard tail fit directly with a primary Gaussian would mix
the tail change with the window and background change. The wider window
has not undergone an independent support or coverage qualification.

\begin{table}[H]\centering\small
\begin{tabular}{llrrrr}
\toprule
Dataset & Shape & $m_0$ (MeV) & $Z_0$ & Local $p_0$ & $\varepsilon^2_{90}$ \\
\midrule
2015 & Gaussian & 93 & 3.52179 & $2.143\times10^{-4}$ & $1.453\times10^{-4}$ \\
2015 & Dilation +10\% & 93 & 3.51668 & $2.185\times10^{-4}$ & $1.457\times10^{-4}$ \\
2015 & Dilation +20\% & 93 & 3.51147 & $2.228\times10^{-4}$ & $1.461\times10^{-4}$ \\
2015 & Dilation +30\% & 93 & 3.50616 & $2.273\times10^{-4}$ & $1.465\times10^{-4}$ \\
2015 & Curvature +10\% & 93 & 3.52029 & $2.155\times10^{-4}$ & $1.454\times10^{-4}$ \\
2015 & Curvature +20\% & 93 & 3.51907 & $2.165\times10^{-4}$ & $1.455\times10^{-4}$ \\
2015 & Curvature +30\% & 93 & 3.51806 & $2.174\times10^{-4}$ & $1.456\times10^{-4}$ \\
2016 & Gaussian & 91 & 3.46714 & $2.630\times10^{-4}$ & $1.388\times10^{-5}$ \\
2016 & Dilation +10\% & 91 & 3.46647 & $2.637\times10^{-4}$ & $1.392\times10^{-5}$ \\
2016 & Dilation +20\% & 91 & 3.46582 & $2.643\times10^{-4}$ & $1.397\times10^{-5}$ \\
2016 & Dilation +30\% & 91 & 3.46513 & $2.650\times10^{-4}$ & $1.401\times10^{-5}$ \\
2016 & Curvature +10\% & 91 & 3.46701 & $2.631\times10^{-4}$ & $1.389\times10^{-5}$ \\
2016 & Curvature +20\% & 91 & 3.46691 & $2.632\times10^{-4}$ & $1.390\times10^{-5}$ \\
2016 & Curvature +30\% & 91 & 3.46682 & $2.633\times10^{-4}$ & $1.391\times10^{-5}$ \\
2021 & Gaussian & 78 & 2.76302 & $2.863\times10^{-3}$ & $6.469\times10^{-6}$ \\
2021 & Dilation +10\% & 78 & 2.76074 & $2.884\times10^{-3}$ & $6.486\times10^{-6}$ \\
2021 & Dilation +20\% & 78 & 2.75828 & $2.905\times10^{-3}$ & $6.504\times10^{-6}$ \\
2021 & Dilation +30\% & 78 & 2.75564 & $2.929\times10^{-3}$ & $6.521\times10^{-6}$ \\
2021 & Curvature +10\% & 78 & 2.76235 & $2.869\times10^{-3}$ & $6.474\times10^{-6}$ \\
2021 & Curvature +20\% & 78 & 2.76180 & $2.874\times10^{-3}$ & $6.478\times10^{-6}$ \\
2021 & Curvature +30\% & 78 & 2.76132 & $2.878\times10^{-3}$ & $6.481\times10^{-6}$ \\
\bottomrule
\end{tabular}

\caption{Guard-window scan minima and their corresponding observed
limits. These rows are separate fixed-shape local calculations, with a
Gaussian guard-window reference.}
\label{tab:guard-peaks}
\end{table}

\clearpage
\subsection{Guard-window 2015 comparison}
Widening the excluded region removes sideband information and adds
explicit tail-sensitive bins to the likelihood. Both effects are included
in the Gaussian guard reference. Near a support edge, the available
external sidebands need not be symmetric; the fixed support and actual
bin selection remain part of the definition.

\studfig{0.98}{0.72}{../figures/guard_2015.pdf}{Observed 2015
$\pm4.5\sig$ fits with external GP training. The 91--100~MeV endpoint
prescription is retained. Tail effects are assessed relative to the
Gaussian fitted with this same wider mask.}

\clearpage
\subsection{Guard-window 2016 comparison}
Broader fitted structure can overlap background directions permitted by
the correlated GP constraint. The profiled likelihood accounts for that
overlap through its nuisance parameters. An observed change in $p_0$ or
the upper limit consequently measures the response of the complete
specified fit, rather than a direct measurement of the detector's signal
tail.

\studfig{0.98}{0.72}{../figures/guard_2016.pdf}{Observed 2016
$\pm4.5\sig$ fits over 39--175~MeV. The final five primary-scan masses
fail the external-sideband geometric requirement for this wider mask.
The GP mean and covariance are held fixed across
signal hypotheses within each mass and window; nuisance shifts are
profiled at every tested amplitude. The archived 2016 state qualification
remains in force.}

\clearpage
\subsection{Guard-window 2021 comparison}
The 2021 guard scan retains the original 10\% exposure and normalization.
It is not an exposure projection. Its broader likelihood region tests
whether the declared outer template changes the pointwise inference
when those bins are used explicitly.

\studfig{0.98}{0.72}{../figures/guard_2021.pdf}{Observed 2021 10\%
$\pm4.5\sig$ fits. Each alternative is compared with the Gaussian under
the same support, external training set and fit window.}

\clearpage
\section{Size and location of the response}
The ratios below compare each alternative with the Gaussian from the
same dataset, mass and window. They quantify changes in the reported
observed limit, not expected sensitivity. The collection of scenarios
is not a confidence band or a nuisance-profiled uncertainty interval.

\begin{table}[H]\centering\footnotesize
\renewcommand{\arraystretch}{0.98}
\begin{tabular}{llrrr}
\toprule
Dataset & Shape & Primary median (\%) & Guard range (\%) & Max. $|\Delta r|$ (guard) \\
\midrule
2015 & Dilation +10\% & 0.168 & -0.030--0.687 & 0.01276 \\
2015 & Dilation +20\% & 0.345 & -0.047--1.378 & 0.02537 \\
2015 & Dilation +30\% & 0.529 & -0.048--2.068 & 0.03773 \\
2015 & Curvature +10\% & 0.049 & -0.007--0.195 & 0.00358 \\
2015 & Curvature +20\% & 0.089 & -0.010--0.352 & 0.00642 \\
2015 & Curvature +30\% & 0.123 & -0.015--0.480 & 0.00871 \\
2016 & Dilation +10\% & 0.168 & 0.065--0.649 & 0.00866 \\
2016 & Dilation +20\% & 0.345 & 0.147--1.336 & 0.01804 \\
2016 & Dilation +30\% & 0.529 & 0.245--2.056 & 0.02796 \\
2016 & Curvature +10\% & 0.049 & 0.025--0.192 & 0.00268 \\
2016 & Curvature +20\% & 0.089 & 0.048--0.351 & 0.00491 \\
2016 & Curvature +30\% & 0.123 & 0.068--0.483 & 0.00678 \\
2021 & Dilation +10\% & 0.168 & 0.087--0.517 & 0.00623 \\
2021 & Dilation +20\% & 0.345 & 0.196--1.032 & 0.01215 \\
2021 & Dilation +30\% & 0.529 & 0.279--1.543 & 0.01772 \\
2021 & Curvature +10\% & 0.049 & 0.023--0.143 & 0.00167 \\
2021 & Curvature +20\% & 0.089 & 0.036--0.257 & 0.00310 \\
2021 & Curvature +30\% & 0.123 & 0.042--0.349 & 0.00439 \\
\bottomrule
\end{tabular}

\caption{Observed percent limit changes relative to the Gaussian in the
same window. Positive values denote larger upper limits.}
\end{table}

\studfig{0.98}{0.34}{../figures/structure_response.pdf}{Response of the
observed inference at the largest declared tail change, $k=1.3$:
percent changes in the limit and changes in the signed likelihood root
are compared with the Gaussian at each mass. Both tail families are
shown, and primary and guard results retain their separate Gaussian
references. Features in a mass scan reuse sideband information and
should not be counted as independent evidence.}

\clearpage
\section{Signal in the GP training region}
Signal tails can enter bins used for background-only GP training.
Changing the local signal template does not correct that contamination.
A wider mask can reduce the training contribution, while also removing
background information.

The diagnostic uses anchors 51 and 92~MeV in 2015, 90 and 92~MeV in
2016, and 78 and 92~MeV in 2021. At each anchor the full injected yield
is fixed to the Gaussian primary observed $A_{90}$, retaining the nominal
conversion. Each window has its own smooth source $n_0$, the full-support
GP prediction from the observed external bins. Reconditioning on the
external bins of $n_0$ gives $b_0,C_0$. The deterministic fitted-bin
expectation is then $n=b_0+A_{\rm true}T_{\rm truth}$.

The clean-training branch fits this expectation with $b_0,C_0$. The
contaminated-training branch adds the same full signal to $n_0$ before
reconditioning, obtains $b_1,C_1$, and fits exactly the same fitted-bin
expectation. Both branches are evaluated with the matching signal
template and with a Gaussian extraction template. Comparing
$\widehat A/A_{\rm true}$ isolates the declared training response and
template mismatch. Clean matched recovery equals one by construction:
centering on $b_0$ makes this a numerical identity, not an independent
background or injection-closure test. Primary and guard sources are
constructed separately, so their difference also changes the reference
continuum; it is not a comparison of windows under one common truth.

\studfig{0.98}{0.42}{../figures/leakage.pdf}{Deterministic signal-leakage
diagnostic: recovered yield divided by injected yield for contaminated
training with matched Gaussian or smooth-dilation extraction. The
horizontal line marks unit recovery. Clean-training, Gaussian-mismatch
and curvature controls are retained in the numerical ledger. The
injected yield is the Gaussian primary observed $A_{90}$ at each anchor.}

No random sampling or coverage test is performed. These conditional
responses do not establish that an observed structure arose from signal
leakage or that an alternative template is the true detector response.

\clearpage
\section{Likelihood and probability construction}
For one dataset, mass, window and shape, the observed expectation is
\begin{equation}
 \lambda_i=b_i+(L\bm\theta)_i+A T_i,\qquad LL^T=C,
\end{equation}
with profiled negative log likelihood, up to a data-only constant,
\begin{equation}
 \ell(A,\bm\theta)=
 \sum_i\left[\lambda_i-n_i+n_i\log\frac{n_i}{\lambda_i}\right]
 +\frac12\bm\theta^T\bm\theta.
 \label{eq:nll}
\end{equation}
The zero-count contribution is $\lambda_i$. The GP count mean $b$ and
covariance $C$ are determined by the external sidebands. Poisson
expectations are constrained to remain physical. The same $b,C,T$ enter
the null, unrestricted and tested-amplitude fits at a fixed coordinate.
The unrestricted amplitude may be negative, enabling a signed diagnostic;
the physical signal strength for limits is nonnegative.

Writing $\widehat A$ for the unrestricted best fit, define
\begin{align}
 r&=\operatorname{sign}(\widehat A)
 \sqrt{2[\ell(0,\widehat{\bm\theta}_0)
               -\ell(\widehat A,\widehat{\bm\theta})]},\\
 Z_0&=\max(r,0),\qquad p_0=\Phi(-Z_0).
\end{align}
Thus a deficit has $Z_0=0$ and $p_0=0.5$ by the inherited one-sided
reference convention. This is distinct from the signed diagnostic
$p_{\mathrm{signed}}=\Phi(-r)$.

For a tested $A\geq0$, the bounded upper-limit statistic is zero if
$A<\widehat A$. Otherwise its denominator uses the unrestricted fit
when $\widehat A\geq0$ and the null fit when $\widehat A<0$.
Let $q$ be that observed statistic and $q_{\mathrm A}$ the explicitly
profiled background-Asimov statistic at the same tested amplitude.
With $a_{\mathrm A}=\sqrt{q_{\mathrm A}}$, the asymptotic tails are
\begin{equation}
 (z_{s+b},z_b)=
 \begin{cases}
 (\sqrt q,\sqrt q-a_{\mathrm A}),&q\leq q_{\mathrm A},\\[2pt]
 \left(\dfrac{q+q_{\mathrm A}}{2a_{\mathrm A}},
       \dfrac{q-q_{\mathrm A}}{2a_{\mathrm A}}\right),&q>q_{\mathrm A},
 \end{cases}
 \qquad
 CL_s(A)=\frac{\Phi(-z_{s+b})}{\Phi(-z_b)}.
\end{equation}
The $q=0$ inclusive-tail convention gives $CL_s=1$. The observed
upper limit solves $CL_s(A_{90})=0.10$. Both likelihood ratios are
profiled explicitly; their sampling tails use the asymptotic
approximation~\cite{Cowan2011}. No Gaussian error-bar substitution is
used for the observed upper limit.

\clearpage
\section{Input provenance, checks and interpretation}
The portable spectra and fitting machinery come from Analysis Note
v5.0.5~\cite{HPSv505}. Native counts, bin edges, resolution, normalization
and archived GP states are preserved, with original ROOT identities in
the export provenance and exact copied-file hashes in the study manifest.

The Gaussian comparison is recomputed rather than copied from the
historical results table. Its primary peak significances are 3.139237,
3.424752 and 2.808645 for 2015, 2016 and 2021, compared with the
published-note values 3.139465, 3.423687 and 2.810290. Across the full
primary scans the largest difference from the inherited table is
$|\Delta Z|=0.12366$, at 83~MeV in 2016; the largest relative limit
difference is about 1.14\%. At 2016/83~MeV the legacy L-BFGS null fit
stopped on relative objective change with maximum gradient about 0.116;
its $Z_0=0.24369$ becomes 0.12003 in the converged fit. Holding the
background, covariance and template fixed reproduces this difference;
the covariance-factor choice changes the converged limit by less than
$6\times10^{-12}$ relative. This numerical issue is distinct from the
2016 kernel-fit caveat. All tail comparisons use the newly converged
Gaussian with identical settings.

The RBF kernel acts in log mass with the inherited log-count noise
prescription. Its coordinates remain fixed across shapes. Effective
radiative fractions are 0.07905, 0.04650 and 0.04770 for the three
datasets, including the inherited penalties. No setting is selected by
an observed limit or probability.

\studtable{../derived/qa_table.tex}{Recorded numerical and artifact
checks. Passing these checks establishes the stated numerical
calculation; it does not establish frequentist coverage or the physical
correctness of a chosen tail model.}

Choosing a preferred shape requires an independent detector-response
model or a calibrated selection procedure. These scenario comparisons
supply no new expected bands, coherent scan calibration, look-elsewhere
correction or global probability.

\begin{thebibliography}{9}
\bibitem{HPSv505}
E.~Peets, M.~Gignac and E.~Hsu,
\emph{HPS Gaussian-Process Resonance Search: Analysis Note v5.0.5},
16 September 2026, internal review draft and archived source package.
\bibitem{Cowan2011}
G.~Cowan, K.~Cranmer, E.~Gross and O.~Vitells,
``Asymptotic formulae for likelihood-based tests of new physics,''
Eur. Phys. J. C \textbf{71}, 1554 (2011),
\href{https://arxiv.org/abs/1007.1727}{arXiv:1007.1727}.
\end{thebibliography}
\end{document}
